\section{Introduction}

Detectors of machine-generated text report accuracy at or above 97\% on the benchmarks the field uses most, and often above 99\%~\cite{Guo2023close,Su2023plus,Zhou2024exploiting}. Read directly, those numbers say the task is close to solved. Read carefully, they say something weaker. A benchmark score measures the separability of a particular corpus, not the capability a reader infers from it.

Separability has more than one source. A detector can reach a high score by modelling how machine-generated language differs from human language. It can also reach it by modelling how one corpus was assembled, punctuated and encoded. Both produce the same headline figure. Individual instances of the second are documented. They include a whitespace convention in HC3~\cite{Tian2023multiscale}, a length confound in the same corpus~\cite{Baidya2026detecting} and prompt-specific collection shortcuts~\cite{Park2024investigating}. What is missing is a way to ask how much, in total, a benchmark's separability owes to surface form rather than content, and to place two benchmarks on the same axis.

Three lines of prior work come close, and each stops short. Shared benchmarks such as RAID~\cite{Dugan2024raid} and M4~\cite{Wang2024m4} make detectors comparable across conditions, but do not ask what any one condition is separable by. Cleaning kits remove a single known cue~\cite{Tian2023multiscale}, but a null after cleaning says nothing unless the detector could read the cue. Partial-input baselines expose artefacts in natural language inference~\cite{Gururangan2018annotation,Poliak2018hypothesis}, but have not been transferred to detection benchmarks, where the artefact is a property of the text rather than a span of it.

This paper proposes that measurement. We compare three arms on identical splits. A \emph{surface-only} model reads 47 orthographic features and never sees a word. A \emph{content-only} model reads text after punctuation, casing and non-ASCII characters are stripped. A \emph{full} model is a fine-tuned transformer on raw text. The first two share a classifier family and are directly comparable. The third is a reference point, not a matched arm. The object of measurement is a corpus, not a detector, and we offer no detector comparison, since a fair one would need matched training data on both corpora. On HC3 the two arms are indistinguishable, 3.20\% against 3.26\% test error, while on DAIGT V2 content is stronger by a factor of 7.9. Splitting each corpus by its own subgroup labels then shows the HC3 tie belongs to the Reddit sub-corpus that is three quarters of it, and reverses on the other domains. The main contributions are the following.
\begin{itemize}
\item A surface-content decomposition, stated formally in Section~\ref{sec:decomp-def} as a reusable measurement over any human-versus-machine corpus, with the length control that stops its two arms sharing a channel.
\item Its application to two benchmarks and their eighteen sub-corpora (Sections~\ref{sec:decomp} and~\ref{sec:subgroup}). This locates the HC3 result in one collection convention in one dominant domain, and shows that removing that convention costs the surface arm ten points.
\item A tokenisation control (Section~\ref{sec:tok}) establishing that the most-discussed cue in HC3 is sufficient in isolation yet unnecessary in practice, since BERT cannot represent it and still reaches 0.9916 weighted F1 there.
\end{itemize}

The rest of this paper is organised as follows. Section~\ref{sec:related} reviews detection on these benchmarks, shortcut learning and robustness. Section~\ref{sec:method} states the problem setting, the corpora, the decomposition and the controls. Section~\ref{sec:results} reports the results. Section~\ref{sec:limits} discusses what they license, states the limits and concludes.
